\documentclass[10pt, aspectratio=169]{beamer}
\usetheme{PaloAlto}
\usecolortheme{seahorse}
\setbeamertemplate{navigation symbols}{}
\setbeamertemplate{caption}[numbered]
%general package
\usepackage[utf8]{inputenc}
\usepackage{array}
\usepackage[english]{babel}
\usepackage{geometry}
\usepackage{tcolorbox}
\usepackage[export]{adjustbox}
\usepackage{graphicx}
\usepackage{xcolor}
\graphicspath{{../img}}
\hypersetup{
    colorlinks=true, 
    linkcolor=red,    % Internal links
    urlcolor=blue     % External links
}
\usepackage{wrapfig}
\usepackage{luacode}
\usepackage{comment}
\usepackage{ulem}
%math package
\usepackage{amsmath}
\usepackage{amsfonts}
%\usepackage{amssymb}
%\usepackage{amsthm}
%\usepackage{slashed}
%\usepackage{tikz-cd}

%font package
%\usepackage{mathrsfs}
%\usepackage{bm}

%misc. package
\usepackage{enumitem}
\usepackage{animate}

\DeclareMathOperator{\xd}{\,d\!}
\DeclareMathOperator{\curl}{curl}
\DeclareMathOperator{\dive}{div}
\newcommand{\e}{{\rm e}}
\newcommand{\norm}[1]{\lVert#1\rVert}
\newcommand{\R}{\mathbb R}
\newcommand{\vF}{\mathbf F}
\newcommand{\vv}{\mathbf v}
\newcommand{\inpr}[1]{\left\langle#1\right\rangle}
\newcommand{\blankwans}[3][.]{\alt<#2>{{\color{red}\uline{#3}}#1}{\hrulefill#1}}

\author[B.H.]{{\Large A Collection of Daily Exercises}\\\vspace{.5em} Author: Benjamin Huang}
\date{}
\title[Skill Builder]{Math Skill Builder}

%\institute[]{\vskip -2em\includegraphics[width = 0.65\textwidth]{}}
%\logo{\includegraphics[width=.12\textwidth]{}}

\begin{document}
\begin{luacode*}
function generateNonzero(min, max)
  local output
  repeat
    output = math.random(min, max)
  until output ~= 0
  return output
end
\end{luacode*}
\frame{\titlepage}
\Large
\begin{frame}
\begin{itemize}
\item
What are the $x$-intercepts that you can see on the graph? \pause
\vspace{\stretch{1}}
\item
Are there any other  $x$-intercepts on the graph outside of the graphing window? How do you know?\pause
\vspace{\stretch{1}}
\item
Can you come up with a polynomial that have no $x$-intercept?
\end{itemize}
\end{frame}

\begin{frame}
\begin{itemize}
\item
In mathematics, a \textbf{product} means \alt<2->{\underline{\color{red}the result of a multiplication}.}{\underline{\phantom{the result of a multiplication}}.} \pause\pause
\vspace{\stretch{1}}
\item
 A \textbf{factor} means \blankwans{4-}{one of the numbers/expressions that form the product}
\end{itemize}
\end{frame}

\begin{frame}
\pause
\begin{itemize}
\item
If the product is zero, then  \alt<3->{\underline{\color{red}one of the factors is zero}.}{\underline{\phantom{one of the factors is zero}}.}\pause\pause
\vspace{\stretch{1}}
\item
If the product is non-zero, then \alt<5->{\underline{\color{red}all of the factors are non-zero}.}{\underline{\phantom{all of the factors are non-zero}}.}\pause \pause
\vspace{\stretch{1}}
\item
If $x_0$ is an $x$-intercept of $f(x)$, then \alt<7->{\underline{\color{red}$f(x_0) = 0$}.}{\underline{\phantom{$f(x_0) = 0$}}.}\pause \pause 
\vspace{\stretch{1}}
\item  
If $y_0$ is a $y$-intercept of $f(x)$, then \alt<9->{\underline{\color{red}$y_0 = f(0)$}.}{\underline{\phantom{$y_0 = f(0)$}}.}
\vspace{\stretch{1}}

\end{itemize}
\end{frame}

\begin{frame}
\begin{itemize}
\item 
The zeros of the function $P(x)$ are \blankwans{2-}{the solutions of the equation $P(x) = 0$, e.g., $1$ and $2$ are the zeros of $x^2 -3x + 2$}\pause \pause 
\vspace{\stretch{1}}
\item
A polynomial is in the \textbf{factored form} if \blankwans{4-}{it's written as a product of first degree polynomial(s), e.g., $(x+2)(x-2)$ is the factored form of $x^2 - 4$}\pause\pause
\vspace{\stretch{1}}
\item
The factored form of a polynomial is good for \blankwans{6-}{relating the zeros of the polynomial immediately, e.g., $-2$ and $3$ are the only zeros of $(x+2)^8(x-3)^5$}
\end{itemize}
\end{frame}


\begin{frame}
\begin{itemize}
\item
The degree of a polynomial is \blankwans{2-}{the highest exponent of the variable, e.g., the degree of $-3 + 2x + 100x^2 - 4x^3$ is 3}\pause\pause
\vspace{\stretch{1}}
\item
A term in a sum is \blankwans{4-}{one of the numbers/expressions that form the sum, e.g., in $3x^2 - 2x + 1$, the terms are $3x^2$, $-2x$ and $1$}\pause\pause
\vspace{\stretch{1}}
\item 
The leading term of a polynomial is \blankwans{6-}{the term of the variable with the highest exponent, e.g., the leading term of $-3 + 2x + 100x^2 - 4x^3$ is $-4x^3$}
\vspace{\stretch{1}}
\end{itemize}

\end{frame}

\begin{frame}
\begin{itemize}
\item
The standard form of a polynomial is \blankwans{2-}{the sum with the terms ordered by exponents from high to low, e.g., $3x^3 - 2x + 4$}
\pause\pause
\vspace{\stretch{1}}
\item
The standard form is good for \blankwans{4-}{identifying the degree, the constant value (a.k.a the $y$-intercept), and implementing the polynomial long division}
\end{itemize}

\end{frame}

\begin{frame}
\begin{itemize}
\item
Before you plot the graph of a function on a graphing calculator, you should \blankwans{2-}{set the graphing window properly}
\pause \pause
\vspace{\stretch{1}}

\item 
To set the graphing window, firstly, you should \blankwans{4-}{find all the $x$-intercepts.}
\end{itemize}
\end{frame}

\end{document}